\documentclass[10pt,a4paper]{amsart}
\usepackage{amsmath,amssymb,amsthm,graphicx}
% Retained preamble for the Stacks Project chapter snapshot in this directory.
% The chapter's own preamble.tex is a copy of the shared Stacks Project preamble; the block
% below is sliced MECHANICALLY, line for line, from it. Nothing was added, removed or
% redefined; the wrapper class, the xr-hyper/hyperref wiring and the externaldocument list
% are omitted because the assembler supplies the document class.
%
% lines taken from preamble.tex: 3, 5, 9, 11, 12, 156, 157, 158, ... (51 lines total)
%
% package coverage restored from the chapter's own preamble (2026-10-13): xy, xy-2cell

\usepackage[all]{xy}
\xyoption{2cell}
\UseAllTwocells

\newenvironment{reference}{\comment}{\endcomment}
\newenvironment{slogan}{\comment}{\endcomment}
\newenvironment{history}{\comment}{\endcomment}
\theoremstyle{plain}
\newtheorem{theorem}[subsection]{Theorem}
\newtheorem{proposition}[subsection]{Proposition}
\newtheorem{lemma}[subsection]{Lemma}
\theoremstyle{definition}
\newtheorem{definition}[subsection]{Definition}
\newtheorem{example}[subsection]{Example}
\newtheorem{exercise}[subsection]{Exercise}
\newtheorem{situation}[subsection]{Situation}
\theoremstyle{remark}
\newtheorem{remark}[subsection]{Remark}
\newtheorem{remarks}[subsection]{Remarks}
\numberwithin{equation}{subsection}
\def\lim{\mathop{\mathrm{lim}}\nolimits}
\def\colim{\mathop{\mathrm{colim}}\nolimits}
\def\Spec{\mathop{\mathrm{Spec}}}
\def\Hom{\mathop{\mathrm{Hom}}\nolimits}
\def\Ext{\mathop{\mathrm{Ext}}\nolimits}
\def\SheafHom{\mathop{\mathcal{H}\!\mathit{om}}\nolimits}
\def\SheafExt{\mathop{\mathcal{E}\!\mathit{xt}}\nolimits}
\def\Sch{\mathit{Sch}}
\def\Mor{\mathop{\mathrm{Mor}}\nolimits}
\def\Ob{\mathop{\mathrm{Ob}}\nolimits}
\def\Sh{\mathop{\mathit{Sh}}\nolimits}
\def\NL{\mathop{N\!L}\nolimits}
\def\CH{\mathop{\mathrm{CH}}\nolimits}
\def\proetale{{pro\text{-}\acute{e}tale}}
\def\etale{{\acute{e}tale}}
\def\QCoh{\mathit{QCoh}}
\def\Ker{\mathop{\mathrm{Ker}}}
\def\Im{\mathop{\mathrm{Im}}}
\def\Coker{\mathop{\mathrm{Coker}}}
\def\Coim{\mathop{\mathrm{Coim}}}
\DeclareMathSymbol{\boxtimes}{\mathbin}{AMSa}{"02}
\def\QCohstack{\mathcal{QC}\!\mathit{oh}}
\def\Cohstack{\mathcal{C}\!\mathit{oh}}
\def\Spacesstack{\mathcal{S}\!\mathit{paces}}
\def\Quotfunctor{\mathrm{Quot}}
\def\Hilbfunctor{\mathrm{Hilb}}
\def\Curvesstack{\mathcal{C}\!\mathit{urves}}
\def\Polarizedstack{\mathcal{P}\!\mathit{olarized}}
\def\Complexesstack{\mathcal{C}\!\mathit{omplexes}}
\def\Pic{\mathop{\mathrm{Pic}}\nolimits}
\def\Picardstack{\mathcal{P}\!\mathit{ic}}
\def\Picardfunctor{\mathrm{Pic}}
\def\Deformationcategory{\mathcal{D}\!\mathit{ef}}
\usepackage{cleveref}
\crefname{originallemma}{original lemma}{original lemmas}
\Crefname{originallemma}{Original lemma}{Original lemmas}
\makeatletter
\newlabel{derived-lemma-representable-homological}{{0149}{source}{}{sourcepointer-1}{}}
\newlabel{derived-lemma-representable-homological@cref}{{[originallemma][0][]0149}{[1][0][]source}{}{}{}}
\makeatother
\crefname{originallemma}{original lemma}{original lemmas}
\Crefname{originallemma}{Original lemma}{Original lemmas}
\makeatletter
\newlabel{lemma-canonical-class-obstruction}{{0DIW}{source}{}{sourcepointer-2}{}}
\newlabel{lemma-canonical-class-obstruction@cref}{{[originallemma][0][]0DIW}{[1][0][]source}{}{}{}}
\makeatother
\crefname{originallemma}{original lemma}{original lemmas}
\Crefname{originallemma}{Original lemma}{Original lemmas}
\makeatletter
\newlabel{lemma-inf-obs-map-defo-complex}{{0DIY}{source}{}{sourcepointer-3}{}}
\newlabel{lemma-inf-obs-map-defo-complex@cref}{{[originallemma][0][]0DIY}{[1][0][]source}{}{}{}}
\makeatother
\crefname{originallemma}{original lemma}{original lemmas}
\Crefname{originallemma}{Original lemma}{Original lemmas}
\makeatletter
\newlabel{lemma-lift-map-complexes}{{0DIX}{source}{}{sourcepointer-4}{}}
\newlabel{lemma-lift-map-complexes@cref}{{[originallemma][0][]0DIX}{[1][0][]source}{}{}{}}
\makeatother
\title{Stacks, Deformation Theory: lemma-first-order-defos-complex}
\author{The Stacks Project Authors}
\date{stacks-project:defos.tex:local-snapshot-20261012; GNU FDL 1.2 or later}
\begin{document}
\maketitle
\noindent\textit{Editorial note.} Statement span is lines 5802-5824 of sources/stacks-defos/pinned-original-source.tex; proof span is lines 5826-5921, the complete author proof (94 body lines).\par\medskip
\noindent\textit{Editorial note.} Navigation stubs: 4 cited-but-not-included labels are supplied as non-load-bearing pointers with their pinned tags where the snapshot tag index has them.\par\medskip
\noindent\textit{Editorial note.} extra\_packages is empty on purpose: the retained preamble already loads xy with options and a second declaration clashes in the assembler wrapper.\par\medskip
\noindent\textit{Editorial note.} Text is retained verbatim, including the author's own typos and placeholders.\par\medskip
\noindent\textit{Editorial note.} Queue index 5 of queue\_defos.json (pb=3834, stmt 5802-5824, proof 5826-5921).\par\medskip
\noindent\textit{Allowed standard prerequisites (editorial).}\begin{itemize}
\item Explicit prerequisite boundary. The cited inputs are prerequisites; the argument in the retained span is not replaced by any cited result.
\end{itemize}\medskip
\ifdefined\hypertarget\hypertarget{sourcepointer-1}{}\fi
\noindent\textit{Source reference.} Original lemma 0149: sources/stacks-defos/pinned-original-source.tex, source lines 5921--5921. Cited by the retained statement or proof but not part of the collected spans; supplied only as a navigation pointer. The argument does not depend on its content.\par\medskip
\ifdefined\hypertarget\hypertarget{sourcepointer-2}{}\fi
\noindent\textit{Source reference.} Original lemma 0DIW: sources/stacks-defos/pinned-original-source.tex, source lines 5575--5588. Cited by the retained statement or proof but not part of the collected spans; supplied only as a navigation pointer. The argument does not depend on its content.\par\medskip
\ifdefined\hypertarget\hypertarget{sourcepointer-3}{}\fi
\noindent\textit{Source reference.} Original lemma 0DIY: sources/stacks-defos/pinned-original-source.tex, source lines 5753--5768. Cited by the retained statement or proof but not part of the collected spans; supplied only as a navigation pointer. The argument does not depend on its content.\par\medskip
\ifdefined\hypertarget\hypertarget{sourcepointer-4}{}\fi
\noindent\textit{Source reference.} Original lemma 0DIX: sources/stacks-defos/pinned-original-source.tex, source lines 5715--5723. Cited by the retained statement or proof but not part of the collected spans; supplied only as a navigation pointer. The argument does not depend on its content.\par\medskip
\begin{lemma}
\label{lemma-first-order-defos-complex}
Let $\mathcal{C}$ be a site. Let $\mathcal{O} \to \mathcal{O}_0$
be a surjection of sheaves of rings whose kernel is an ideal sheaf
$\mathcal{I}$ of square zero. Let $K_0 \in D^-(\mathcal{O})$.
A lift of $K_0$ is a pair $(K, \alpha_0)$ consisting of an object
$K$ in $D^-(\mathcal{O})$ and an isomorphism
$\alpha_0 : K \otimes_\mathcal{O}^\mathbf{L} \mathcal{O}_0 \to K_0$
in $D(\mathcal{O}_0)$.
\begin{enumerate}
\item Given a lift $(K, \alpha)$ the group of automorphism of the pair
is canonically the cokernel of a map
$$
\Ext^{-1}_{\mathcal{O}_0}(K_0, K_0)
\longrightarrow
\Hom_{\mathcal{O}_0}(K_0, K_0 \otimes_{\mathcal{O}_0}^\mathbf{L} \mathcal{I})
$$
\item If there is a lift, then the set of isomorphism classes of lifts
is principal homogenenous under
$\Ext^1_{\mathcal{O}_0}(K_0,
K_0 \otimes_{\mathcal{O}_0}^\mathbf{L} \mathcal{I})$.
\end{enumerate}
\end{lemma}

\begin{proof}
An automorphism of $(K, \alpha)$ is a map $\varphi : K \to K$
in $D(\mathcal{O})$ with
$\varphi \otimes_\mathcal{O} \text{id}_{\mathcal{O}_0} = \text{id}$.
This is the same thing as saying that
$$
K \xrightarrow{\varphi - \text{id}} K \to
K \otimes_\mathcal{O}^\mathbf{L} \mathcal{O}_0
$$
is zero. We conclude the group of automorphisms is
the cokernel of a map
$$
\Hom_\mathcal{O}(K, K_0[-1])
\longrightarrow
\Hom_\mathcal{O}(K, K_0 \otimes_{\mathcal{O}_0}^\mathbf{L} \mathcal{I})
$$
by the distinguished triangle
$$
K \otimes_\mathcal{O}^\mathbf{L} \mathcal{I} \to
K \to
K \otimes_\mathcal{O}^\mathbf{L} \mathcal{O}_0 \to
(K \otimes_\mathcal{O}^\mathbf{L} \mathcal{I})[1]
$$
in $D(\mathcal{O})$ and 
Derived Categories, Lemma \ref{derived-lemma-representable-homological}.
To translate into the groups in the lemma use adjunction
of the restriction functor $D(\mathcal{O}_0) \to D(\mathcal{O})$ and
$- \otimes_\mathcal{O} \mathcal{O}_0 : D(\mathcal{O}) \to D(\mathcal{O}_0)$.
This proves (1).

\medskip\noindent
Proof of (2).
Assume that $K_0 = K \otimes_\mathcal{O}^\mathbf{L} \mathcal{O}_0$
in $D(\mathcal{O})$. By Lemma \ref{lemma-inf-obs-map-defo-complex}
the map sending a lift $(K', \alpha_0)$ to the obstruction $o(\alpha_0)$
to lifting $\alpha_0$ defines a canonical injective map
from the set of isomomorphism classes of pairs
to $\Ext^1_{\mathcal{O}_0}(K_0,
K_0 \otimes_{\mathcal{O}_0}^\mathbf{L} \mathcal{I})$.
To finish the proof we show that it is surjective.
Pick $\xi : K_0 \to (K_0 \otimes_{\mathcal{O}_0}^\mathbf{L} \mathcal{I})[1]$
in the $\Ext^1$ of the lemma.
Choose a bounded above complex $\mathcal{F}^\bullet$
of flat $\mathcal{O}$-modules representing $K$.
The map $\xi$ can be represented as $t \circ s^{-1}$
where $s : \mathcal{K}_0^\bullet \to \mathcal{F}_0^\bullet$
is a quasi-isomorphism and
$t : \mathcal{K}_0^\bullet \to
\mathcal{F}_0^\bullet \otimes_{\mathcal{O}_0} \mathcal{I}[1]$
is a map of complexes.
By Lemma \ref{lemma-lift-map-complexes}
we can assume there exists a quasi-isomorphism
$\mathcal{G}^\bullet \to \mathcal{F}^\bullet$
of complexes of $\mathcal{O}$-modules
such that $\mathcal{G}_0^\bullet \to \mathcal{F}_0^\bullet$
factors through $s$ up to homotopy.
We may and do replace $\mathcal{G}^\bullet$ by a bounded
above complex of flat $\mathcal{O}$-modules (by picking a qis
from such to $\mathcal{G}^\bullet$ and replacing).
Then we see that $\xi$ is represented by
a map of complexes
$t : \mathcal{G}_0^\bullet \to
\mathcal{F}_0^\bullet \otimes_{\mathcal{O}_0} \mathcal{I}[1]$
and the quasi-isomorphism $\mathcal{G}_0^\bullet \to \mathcal{F}_0^\bullet$.
Set
$$
\mathcal{H}^n = \mathcal{F}^n \times_{\mathcal{F}_0^n} \mathcal{G}_0^n
$$
with differentials
$$
\mathcal{H}^n \to \mathcal{H}^{n + 1},\quad
(f^n, g_0^n) \mapsto (d(f^n) + t(g_0^n), d(g_0^n))
$$
This makes sense as
$\mathcal{F}_0^{n + 1} \otimes_{\mathcal{O}_0} \mathcal{I} =
\mathcal{F}^{n + 1} \otimes_\mathcal{O} \mathcal{I} =
\mathcal{I}\mathcal{F}^{n + 1} \subset \mathcal{F}^{n + 1}$.
We omit the computation that shows that $\mathcal{H}^\bullet$
is a complex of $\mathcal{O}$-modules. By construction there is
a short exact sequence
$$
0 \to \mathcal{F}_0^\bullet \otimes_{\mathcal{O}_0} \mathcal{I} \to
\mathcal{H}^\bullet \to \mathcal{G}_0^\bullet \to 0
$$
of complexes of $\mathcal{O}$-modules.
Exactly as in the proof of Lemma \ref{lemma-canonical-class-obstruction}
one shows that this sequence induces an isomorphism
$\alpha_0 :
\mathcal{H}^\bullet \otimes_\mathcal{O}^\mathbf{L} \mathcal{O}_0 \to
\mathcal{G}_0^\bullet$ in $D(\mathcal{O}_0)$.
In other words, we have produced a pair $(\mathcal{H}^\bullet, \alpha_0)$.
We omit the verification that $o(\alpha_0) = \xi$; hint: $o(\alpha_0)$
can be computed explicitly in this case as we have maps
$\mathcal{H}^n \to \mathcal{F}^n$ (not compatible with differentials)
lifting the components of $\alpha_0$. This finishes the proof.
\end{proof}

\end{document}
